\documentclass{situationist-preprint}

\title{Religious Language After Verificationism}
\msid{REL-18-a}
\msfamily{REL-18}
\msdiscipline{Philosophy of Religion}
\mstype{Principal manuscript}
\mskeywords{religious language; verificationism; falsification; truth-aptness; compositionality; noncognitivism; analogy}

\begin{document}
\maketitle

\begin{abstract}
Logical positivism held that a sentence is meaningful only if it is analytic or empirically verifiable, and concluded that sentences such as \enquote{God created the world} are neither true nor false. Flew's falsification challenge and Braithwaite's noncognitivism kept the conclusion alive after the criterion had been abandoned. We argue that religious statements are truth-apt, even though most of their truth-conditions cannot be verified. We first show, in the familiar way, that no verifiability criterion separates religious sentences from scientific and historical ones: the criteria either admit everything or exclude much that is plainly meaningful. We then give a positive argument. Truth-aptness is fixed by compositional structure and by the meaning of constituents, and religious sentences are built in the ordinary ways from expressions that have meaning in use, embed in negation, conditionals and attitude reports, and stand in inferential relations recognized by believers and unbelievers alike. That their truth-conditions mostly outrun verification is a feature they share with claims about the remote past and about regions beyond any observer's horizon. We identify which religious claims are partly open to evidence and which are not.
\end{abstract}

\section{The problem}

\citet{Carnap1932} argued that the sentences of metaphysics are pseudo-statements, since their words have no criteria of application or the sentences no conditions of verification. \citet[ch.~6]{Ayer1936} applied the verdict to theology: since no empirical proposition follows from \enquote{God exists}, together with whatever auxiliary premises, that does not follow from the auxiliaries alone, the sentence is literally senseless. Ayer drew the consequence that the atheist and the agnostic are as much in error as the theist, since \enquote{God does not exist} is senseless too.

The verifiability criterion did not survive. \citet{Church1949} showed that the formulation of Ayer's second edition admits any sentence whatever, and \citet{Hempel1950} reviewed the successive formulations and concluded that none draws the line the positivists wanted. But the challenge to religious language survived its criterion. \citet{Flew1955} told the parable of the gardener: two explorers find a clearing, one says a gardener tends it, and as each test fails, the believer qualifies his claim until the gardener is invisible, intangible and undetectable; Flew asked what difference remains between such a gardener and no gardener at all. \citet{Braithwaite1955} proposed that religious assertions express the intention to follow a way of life, associated with stories that need not be believed. Readers of Wittgenstein's lectures on religious belief have found there a view on which religious utterances do not contradict the unbeliever's but express a different picture regulating a life \citep{Wittgenstein1966}. \citet{Hick1960} replied to the verificationist on his own terms: \enquote{God exists} is verifiable in principle, since a person who survived death would find it confirmed in experience. \citet{Swinburne1977} argued that the coherence of theism is a question about whether its claims are consistent, and that statements may be meaningful though unverifiable.

We argue that religious statements are truth-apt, and that this does not depend on Hick's eschatological verification. Their truth-conditions mostly cannot be verified, and that is not a defect of meaning.

\section{Framework}

\begin{definition}[Truth-aptness]\label{d:apt}
A declarative sentence is \emph{truth-apt} if, as used, it expresses a proposition, and so has conditions under which it is true and conditions under which it is false.
\end{definition}

\begin{definition}[Verifiability]\label{d:verif}
A sentence $S$ is \emph{directly verifiable} if it is an observation sentence or entails one. It is \emph{weakly verifiable} (Ayer's first criterion) if there are auxiliary premises $A$ such that $S$ and $A$ together entail an observation sentence that $A$ alone does not entail.
\end{definition}

\begin{principle}[Compositionality]\label{p:comp}
The truth-conditions of a complex sentence are determined by the contributions of its constituents and by its mode of composition.
\end{principle}

\begin{principle}[Meaning in use]\label{p:use}
An expression has a meaning if it has a use governed by rules of application and of inference that competent speakers share, and its contribution to truth-conditions is fixed by that use.
\end{principle}

Principle~\ref{p:use} is the account of LNG-01, on which meaning is use constrained by truth-conditions. It is neutral between theories of reference and does not require that each expression be applicable on the basis of observation; theoretical terms in science are applied through their inferential connections to others, and on the realism of SCI-02 they nonetheless refer.

\section{The argument}

\begin{lemma}[Collapse]\label{l:collapse}
Every sentence is weakly verifiable.
\end{lemma}

\begin{proof}
Let $N$ be any sentence and $O$ any observation sentence not entailed by $N$ alone. Take as auxiliary the conditional $N \rightarrow O$. Then $N$ together with $N \rightarrow O$ entails $O$. The auxiliary alone does not entail $O$, since it is true wherever $N$ is false, including where $O$ is false. By Definition~\ref{d:verif}, $N$ is weakly verifiable.
\end{proof}

Lemma~\ref{l:collapse} is Ayer's own difficulty, which his second edition tried to meet by restricting the auxiliaries; Church's construction shows that the restricted criterion admits any sentence too. Strengthening the criterion to direct verifiability excludes universal laws, which entail no observation sentence by themselves, and claims about the remote past. Since every criterion so far proposed either admits theology or excludes science and history \citep{Hempel1950}, the verificationist argument against religious language has no premise.

\begin{lemma}[Unverifiable truth-apt sentences]\label{l:unverif}
There are truth-apt sentences whose truth-conditions cannot be verified.
\end{lemma}

\begin{proof}
\enquote{The number of hairs on Caesar's head at his death was even} is truth-apt: it is built from expressions with ordinary use, its negation is truth-apt, and exactly one of the two is true. Its truth-conditions cannot be verified, since the evidence has perished. The same holds of claims about events beyond any observer's cosmological horizon.
\end{proof}

\begin{theorem}\label{t:main}
Religious statements such as \enquote{God created the world}, \enquote{God loves every person} and \enquote{the dead will be raised} are truth-apt.
\end{theorem}

\begin{proof}
Each is built by the ordinary modes of composition, predication and quantification, from expressions with meaning in use. \enquote{Created}, \enquote{loves}, \enquote{world}, \enquote{person}, \enquote{dead} and \enquote{raised} have ordinary uses, and their extension to God and to the dead is governed by rules that believers and unbelievers share: that what created the world exists, that whoever loves a person wills her good, that whoever is raised exists after having died. \enquote{God} is used as a name whose reference is fixed by a description, such as \enquote{the maximally great being} of REL-01, or \enquote{the creator of the world}. These uses satisfy Principle~\ref{p:use}. By Principle~\ref{p:comp}, the sentences have truth-conditions. That they express propositions is confirmed by their logical behavior: they embed under negation, in conditionals and in attitude reports (\enquote{if God created the world, the world depends on something}, \enquote{she doubts that God loves her}); they stand in inferential relations that both parties to a dispute acknowledge; and the unbeliever who says \enquote{God did not create the world} contradicts the believer, as the believer takes her to. By Definition~\ref{d:apt}, they are truth-apt. Lemma~\ref{l:unverif} shows that the unverifiability of their truth-conditions, which we examine next, does not count against this.
\end{proof}

\begin{proposition}\label{p:most}
Most of the truth-conditions of religious statements cannot be verified.
\end{proposition}

\begin{proof}
We consider the main classes. Claims about God's necessary existence, timelessness and inner life, such as REL-19's thesis that God is timeless and the doctrine of the Trinity whose consistency REL-13 establishes, have truth-conditions that concern no observable state. Claims about life after death have truth-conditions that, by REL-08, no evidence available before death could confirm. Claims about God's reasons, such as that an outcome was granted in response to a prayer, need not differ in any observable outcome from their negations, as REL-12 argues. Claims about God's creation of the world concern a relation between the world and something outside it, and observation of the world does not settle whether it holds; the cosmological considerations of REL-01 leave the causal role of their object undetermined. Some religious claims are partly open to evidence: the occurrence of reported miracles may be established by testimony (REL-07), and religious experience is prima facie evidence, though of low weight absent cross-checking (REL-21). These bear on a small part of the truth-conditions of the claims they support.
\end{proof}

\section{Predicates applied to God}

A residual worry concerns the predicates. If \enquote{loves} and \enquote{knows} apply to God only in a sense remote from their ordinary one, as the doctrine of analogy holds \citep[I, q.~13]{AquinasST}, then perhaps the sentences in which they occur have no determinate truth-conditions. Principle~\ref{p:use} answers the worry in part. A predicate applied analogically keeps the inferential rules that are not tied to the conditions of its ordinary application: whoever loves another wills her good, whoever knows that $p$ is not mistaken about $p$. These rules fix a contribution to truth-conditions, though not a complete one, and incomplete grasp of a predicate's application is compatible with the sentence's being true or false, as with theoretical terms whose full application is unknown. What the doctrine of analogy shows is that our grasp of the truth-conditions of religious statements is partial. Partial grasp of truth-conditions is not absence of truth-conditions.

\section{Objections and replies}

\begin{objection}[Falsification]
A statement that the believer will not allow anything to count against has no content: it is compatible with every state of affairs. Religious statements, as believers hold them, are of this kind \citep{Flew1955}.
\end{objection}

\begin{reply}
The objection confuses a believer's dispositions with the statement's truth-conditions. That a believer will not give up \enquote{God loves us} in the face of suffering shows how firmly she holds it, not that the statement is compatible with every state of affairs; it is incompatible, for instance, with there being no God. A statement can be incompatible with many states of affairs none of which anyone can observe to obtain. The companion manuscript REL-18-b revisits Flew's parable of the gardener.
\end{reply}

\begin{objection}[Noncognitivism]
Religious utterances express attitudes and commitments to a way of life, not beliefs \citep{Braithwaite1955}. They are no more truth-apt than resolutions.
\end{objection}

\begin{reply}
The practice tells against the description. Believers argue for the existence of God, disagree with atheists, and change their minds on evidence; they treat \enquote{God exists} as something that could be false. Moreover, noncognitivism does not deliver the conclusion it seeks: MTE-03 shows how expressivist contents compose in embedded contexts for moral sentences, and on the deflationism of LNG-10, sentences that are disciplined, embeddable and used assertorically are thereby truth-apt in the minimal sense. Whether deflationism about truth yields minimal truth-aptness so cheaply is disputed; \citet{JacksonOppySmith1994} argue that truth-aptness requires that a sentence be used to express belief. On either view the noncognitivist loses: if minimalism is right, religious sentences are truth-apt because they are disciplined and embeddable; if Jackson, Oppy and Smith are right, they are truth-apt because, as the practice of argument shows, they are used to express beliefs.
\end{reply}

\begin{objection}[Triviality]
If truth-aptness is minimal, Theorem~\ref{t:main} is trivial and says nothing about whether religious statements describe anything.
\end{objection}

\begin{reply}
The argument of Theorem~\ref{t:main} establishes more than minimal truth-aptness: it assigns religious statements truth-conditions through the compositional contribution of expressions whose use connects them to the rest of language, including existential claims whose truth requires that something exist. That is the sense in which \enquote{God created the world} and its negation describe different ways for reality to be. The deflationary point in the previous reply is offered only to show that the noncognitivist cannot avoid truth-aptness even on her own terms.
\end{reply}

\begin{objection}[Eschatological verification suffices]
Hick showed that religious statements are verifiable in principle. The verificationist can be satisfied without abandoning his criterion \citep{Hick1960}.
\end{objection}

\begin{reply}
Hick's argument accepts a criterion that Lemma~\ref{l:collapse} and the history of the criterion show to be untenable, and it covers few religious statements: experience after death would not verify that God is timeless or triune. It also overstates what post-mortem experience could show. REL-08 argues that, under a pattern theory of persons, even the survivor's own experience would not discriminate survival from its matched rivals. Theorem~\ref{t:main} does not need Hick's argument.
\end{reply}

\section{Conclusion}

Religious statements are truth-apt. The verifiability criterion that would have denied this cannot be stated without admitting everything or excluding much that is plainly meaningful. Religious sentences are built in the ordinary ways from expressions with meaning in use, behave logically as truth-apt sentences do, and are contradicted by their denials. Most of their truth-conditions cannot be verified, which they share with claims about the vanished past; some are partly open to testimony and experience. What follows is not that religious language is empty, but that most of what it says lies beyond what evidence can settle.

\begin{thebibliography}{99}\small
\bibitem[Aquinas(ST)]{AquinasST} Aquinas, T. \emph{Summa Theologiae}. Cited by part, question and article.
\bibitem[Ayer(1936)]{Ayer1936} Ayer, A.~J. (1936). \emph{Language, Truth and Logic}. Gollancz. 2nd edn 1946.
\bibitem[Braithwaite(1955)]{Braithwaite1955} Braithwaite, R.~B. (1955). \emph{An Empiricist's View of the Nature of Religious Belief}. Cambridge University Press.
\bibitem[Carnap(1932)]{Carnap1932} Carnap, R. (1932). \"Uberwindung der Metaphysik durch logische Analyse der Sprache. \emph{Erkenntnis}, 2, 219--241.
\bibitem[Church(1949)]{Church1949} Church, A. (1949). Review of A.~J. Ayer, \emph{Language, Truth and Logic}, 2nd edn. \emph{The Journal of Symbolic Logic}, 14(1), 52--53.
\bibitem[Flew(1955)]{Flew1955} Flew, A. (1955). Theology and falsification. In A. Flew and A. MacIntyre (eds), \emph{New Essays in Philosophical Theology}. SCM Press.
\bibitem[Hempel(1950)]{Hempel1950} Hempel, C.~G. (1950). Problems and changes in the empiricist criterion of meaning. \emph{Revue Internationale de Philosophie}, 4(11), 41--63.
\bibitem[Hick(1960)]{Hick1960} Hick, J. (1960). Theology and verification. \emph{Theology Today}, 17(1), 12--31.
\bibitem[Jackson et~al.(1994)]{JacksonOppySmith1994} Jackson, F., Oppy, G. and Smith, M. (1994). Minimalism and truth aptness. \emph{Mind}, 103(411), 287--302.
\bibitem[Swinburne(1977)]{Swinburne1977} Swinburne, R. (1977). \emph{The Coherence of Theism}. Clarendon Press.
\bibitem[Wittgenstein(1966)]{Wittgenstein1966} Wittgenstein, L. (1966). \emph{Lectures and Conversations on Aesthetics, Psychology and Religious Belief}, ed. C. Barrett. Blackwell.
\end{thebibliography}

\end{document}
